% GENERATED FILE. Edit canonical lessons, then run npm run generate:books.
% canonical-source-hash: fnv1a64:a525842ad9972698
% canonical-lessons: BN-C409-elaka, BN-C409-bosti, BN-C409-koloni, BN-C409-basinda, BN-C409-adda-deoya

\chapter{Areas and Localities}
\label{ch:bn-areas-and-localities}
\hlchaptermodality{409}

\begin{chapteropening}
\textbf{By the end of this chapter:} \emph{I can describe an area, a colony and a slum, the people who live there and their evening talk.}

Complete the last lesson of chapter 409: I can describe an area, a colony and a slum, the people who live there and their evening talk.
\end{chapteropening}

\section[\texorpdfstring{\bn{এলাকা} (elākā)}{এলাকা (elākā)}]{\bn{এলাকা} (elākā) — an area; a locality}

\label{lesson:BN-C409-elaka}



\begin{quote}
Before the new one: say the Bengali for a security deposit, then the Bengali for to reside.
\end{quote}

\subsection*{You'll want to know: \bn{এলাকা}}
\textbf{\bn{এলাকা}} — \emph{elākā} — \textquotedblleft{}an area; a locality\textquotedblright{}.

\begin{quote}
\bn{এই} \bn{এলাকায়} \bn{অনেক} \bn{দোকান} \bn{আছে।}
\end{quote}

\begin{quote}
\emph{ei elākāy ônek dokān āchhe.}
\end{quote}

\begin{quote}
There are many shops in this area.
\end{quote}

\begin{cousinweb}
\textbf{\bn{এলাকা}} (\emph{elākā}) is usually traced to Arabic \emph{ʿalāqa}, \textquotedblleft{}a connection, a holding\textquotedblright{}, which in Persian came to mean a district.
\end{cousinweb}

\begin{grammarlens}[title={What you've built}]
One more word for getting around.
\end{grammarlens}

\subsection*{Your turn}
\begin{itemize}
\raggedright
  \item \emph{Say it:} \emph{elākā}
  \item \emph{Say it:} \emph{elākā}, once more
  \item \emph{Say it:} \emph{ei elākāy ônek dokān āchhe.} — There are many shops in this area.
  \item \emph{Recall:} say the Bengali for a security deposit, then the Bengali for to reside, then say \emph{elākā} again
\end{itemize}

\begin{tcolorbox}[breakable,colback=teal!4,colframe=teal!35!black,title={Before you move on}]
What does \bn{এলাকা} mean? (\textquotedblleft{}An area; a locality\textquotedblright{}.) Say it once more.
\end{tcolorbox}
\section[\texorpdfstring{\bn{বস্তি} (bôsti)}{বস্তি (bôsti)}]{\bn{বস্তি} (bôsti) — a slum; a poor settlement}

\label{lesson:BN-C409-bosti}



\begin{quote}
Before the new one: say the Bengali for to reside, then the Bengali for an area.
\end{quote}

\subsection*{You'll want to know: \bn{বস্তি}}
\textbf{\bn{বস্তি}} — \emph{bôsti} — \textquotedblleft{}a slum; a poor settlement\textquotedblright{}.

\begin{quote}
\bn{স্টেশনের} \bn{পাশে} \bn{একটা} \bn{বস্তি} \bn{আছে।}
\end{quote}

\begin{quote}
\emph{sṭeshôner pāshe ekṭā bôsti āchhe.}
\end{quote}

\begin{quote}
There is a slum beside the station.
\end{quote}

\begin{cousinweb}
\textbf{\bn{বস্তি}} (\emph{bôsti}) probably goes back to Sanskrit \emph{vasati}, \textquotedblleft{}a dwelling, a staying\textquotedblright{} (Monier-Williams); in Bengali the word has narrowed to a crowded settlement of the poor.
\end{cousinweb}

\begin{grammarlens}[title={What you've built}]
One more word for getting around.
\end{grammarlens}

\subsection*{Your turn}
\begin{itemize}
\raggedright
  \item \emph{Say it:} \emph{bôsti}
  \item \emph{Say it:} \emph{bôsti}, once more
  \item \emph{Say it:} \emph{sṭeshôner pāshe ekṭā bôsti āchhe.} — There is a slum beside the station.
  \item \emph{Recall:} say the Bengali for to reside, then the Bengali for an area, then say \emph{bôsti} again
\end{itemize}

\begin{tcolorbox}[breakable,colback=teal!4,colframe=teal!35!black,title={Before you move on}]
What does \bn{বস্তি} mean? (\textquotedblleft{}A slum; a poor settlement\textquotedblright{}.) Say it once more.
\end{tcolorbox}
\section[\texorpdfstring{\bn{কলোনি} (kôloni)}{কলোনি (kôloni)}]{\bn{কলোনি} (kôloni) — a housing colony; an estate}

\label{lesson:BN-C409-koloni}



\begin{quote}
Before the new one: say the Bengali for an area, then the Bengali for a slum.
\end{quote}

\subsection*{You'll want to know: \bn{কলোনি}}
\textbf{\bn{কলোনি}} — \emph{kôloni} — \textquotedblleft{}a housing colony; an estate\textquotedblright{}.

\begin{quote}
\bn{মীরার} \bn{মামা} \bn{রেল} \bn{কলোনিতে} \bn{থাকেন।}
\end{quote}

\begin{quote}
\emph{mīrār māmā rel kôlonite thāken.}
\end{quote}

\begin{quote}
Mira's uncle lives in the railway colony.
\end{quote}

\begin{cousinweb}
\textbf{\bn{কলোনি}} (\emph{kôloni}) is borrowed from English \emph{colony}, used in South Asia for a planned group of houses, often for one employer's staff.
\end{cousinweb}

\begin{grammarlens}[title={What you've built}]
One more word for getting around.
\end{grammarlens}

\subsection*{Your turn}
\begin{itemize}
\raggedright
  \item \emph{Say it:} \emph{kôloni}
  \item \emph{Say it:} \emph{kôloni}, once more
  \item \emph{Say it:} \emph{mīrār māmā rel kôlonite thāken.} — Mira's uncle lives in the railway colony.
  \item \emph{Recall:} say the Bengali for an area, then the Bengali for a slum, then say \emph{kôloni} again
\end{itemize}

\begin{tcolorbox}[breakable,colback=teal!4,colframe=teal!35!black,title={Before you move on}]
What does \bn{কলোনি} mean? (\textquotedblleft{}A housing colony; an estate\textquotedblright{}.) Say it once more.
\end{tcolorbox}
\section[\texorpdfstring{\bn{বাসিন্দা} (bāsindā)}{বাসিন্দা (bāsindā)}]{\bn{বাসিন্দা} (bāsindā) — a resident; an inhabitant}

\label{lesson:BN-C409-basinda}



\begin{quote}
Before the new one: say the Bengali for a slum, then the Bengali for a housing colony.
\end{quote}

\subsection*{You'll want to know: \bn{বাসিন্দা}}
\textbf{\bn{বাসিন্দা}} — \emph{bāsindā} — \textquotedblleft{}a resident; an inhabitant\textquotedblright{}.

\begin{quote}
\bn{পাড়ার} \bn{বাসিন্দারা} \bn{রবিবার} \bn{রাস্তা} \bn{পরিষ্কার} \bn{করলেন।}
\end{quote}

\begin{quote}
\emph{pāṛār bāsindārā rôbibār rāstā pôrishkār kôrlen.}
\end{quote}

\begin{quote}
The residents of the neighbourhood cleaned the street on Sunday.
\end{quote}

\begin{cousinweb}
\textbf{\bn{বাসিন্দা}} (\emph{bāsindā}) is usually traced to Persian \emph{bāshinda}, \textquotedblleft{}one who dwells\textquotedblright{}, from \emph{bāshīdan}, to be, to dwell.
\end{cousinweb}

\begin{grammarlens}[title={What you've built}]
One more word for getting around.
\end{grammarlens}

\subsection*{Your turn}
\begin{itemize}
\raggedright
  \item \emph{Say it:} \emph{bāsindā}
  \item \emph{Say it:} \emph{bāsindā}, once more
  \item \emph{Say it:} \emph{pāṛār bāsindārā rôbibār rāstā pôrishkār kôrlen.} — The residents of the neighbourhood cleaned the street on Sunday.
  \item \emph{Recall:} say the Bengali for a slum, then the Bengali for a housing colony, then say \emph{bāsindā} again
\end{itemize}

\begin{tcolorbox}[breakable,colback=teal!4,colframe=teal!35!black,title={Before you move on}]
What does \bn{বাসিন্দা} mean? (\textquotedblleft{}A resident; an inhabitant\textquotedblright{}.) Say it once more.
\end{tcolorbox}
\section[\texorpdfstring{\bn{আড্ডা} \bn{দেওয়া} (āḍḍā deoyā)}{আড্ডা দেওয়া (āḍḍā deoyā)}]{\bn{আড্ডা} \bn{দেওয়া} (āḍḍā deoyā) — to chat in a group; to hang out}

\label{lesson:BN-C409-adda-deoya}



\begin{quote}
Before the new one: say the Bengali for a housing colony, then the Bengali for a resident.
\end{quote}

\subsection*{You'll want to know: \bn{আড্ডা} \bn{দেওয়া}}
\textbf{\bn{আড্ডা} \bn{দেওয়া}} — \emph{āḍḍā deoyā} — \textquotedblleft{}to chat in a group; to hang out\textquotedblright{}.

\begin{quote}
\bn{সন্ধ্যায়} \bn{পাড়ার} \bn{ছেলেরা} \bn{চায়ের} \bn{দোকানে} \bn{আড্ডা} \bn{দেয়।}
\end{quote}

\begin{quote}
\emph{sôndhyāy pāṛār chhelerā chāyer dokāne āḍḍā dey.}
\end{quote}

\begin{quote}
In the evenings the local young men hang out at the tea stall.
\end{quote}

Past (he/she …): \bn{তিনি} \bn{বন্ধুদের} \bn{সঙ্গে} \bn{আড্ডা} \bn{দিলেন}, \emph{tini bôndhuder sôngge āḍḍā dilen}.

\bn{আড্ডা} alone is the gathering itself, a well-loved Bengali habit: long, easy talk over tea about everything and nothing.

\begin{cousinweb}
\textbf{\bn{আড্ডা} \bn{দেওয়া}} (\emph{āḍḍā deoyā}) joins \bn{আড্ডা}, \textquotedblleft{}a gathering place, a long chat\textquotedblright{}, to \bn{দেওয়া}, to give; the origin of \bn{আড্ডা} is not settled.
\end{cousinweb}

\begin{grammarlens}[title={What you've built}]
One more word for getting around.
\end{grammarlens}

\subsection*{Your turn}
\begin{itemize}
\raggedright
  \item \emph{Say it:} \emph{āḍḍā deoyā}
  \item \emph{Say it:} \emph{āḍḍā deoyā}, once more
  \item \emph{Say it:} \emph{sôndhyāy pāṛār chhelerā chāyer dokāne āḍḍā dey.} — In the evenings the local young men hang out at the tea stall.
  \item \emph{Recall:} say the Bengali for a housing colony, then the Bengali for a resident, then say \emph{āḍḍā deoyā} again
\end{itemize}

\begin{tcolorbox}[breakable,colback=teal!4,colframe=teal!35!black,title={Before you move on}]
What does \bn{আড্ডা} \bn{দেওয়া} mean? (\textquotedblleft{}To chat in a group; to hang out\textquotedblright{}.) Say it once more.
\end{tcolorbox}
